\documentclass{exam}

\newif\ifA
\newif\ifkey
\newif\ifprintselected   % required by exam_config's \ansbox

%%%%%%%%%%%%%%%%%%%%%%%%%%%%%%%%%%%%%%%%%%%%%%%%%%%%%%
%%% SET THESE IF VARIABLES
%%%%%%%%%%%%%%%%%%%%%%%%%%%%%%%%%%%%%%%%%%%%%%%%%%%%%%
% Lab name -- labs are named by the lesson they belong to, never by a
% sequential lab number (Neil, 2026-10-12): this is the L20 Lab.
\def\labLesson{20}
%
% The L20 Lab sheet (objective 3.2): Part A of the former L23 Array Factor and
% Pattern Lab sheet, split off 2026-10-13 when each lesson took its own lab.

% is this the key or not?
\ifdefined\iskey
	\ifnum\iskey=1
		\keytrue
	\else
		\keyfalse
	\fi
\else
	%% Pick one manually
	\keyfalse % if this is NOT the key
%	\keytrue     % if this IS the key
\fi

% set up the header and footer - change for every term
\ifdefined\thisTermIn
	\def\thisterm{\thisTermIn}
\else
	\def\thisterm{Fall 2026}
\fi
%
\def\thisclass{ECE 444}

%%%%%%%%%%%%%%%%%%%%%%%%%%%%%%%%%%%%%%%%%%%%%%%%%%%%%%
%%% END SET THESE IF VARIABLES
%%%%%%%%%%%%%%%%%%%%%%%%%%%%%%%%%%%%%%%%%%%%%%%%%%%%%%
% packages and shortcuts
\input{myPackages_gr}
\usepackage[hidelinks, linkcolor=red]{hyperref}
\input{myShortcuts}
\setlength{\parindent}{0pt}
\setlength{\parskip}{6pt}

%% exam class customization
\input{exam_config}
% ECE 444 document fonts.
%
% Body text is Barlow (SIL OFL; the faces and the license are in
% latex/fonts/Barlow/). Barlow is not in TeX Live, so it is loaded as an
% OpenType face through fontspec -- which means the document must be built
% with lualatex, not pdflatex. build_practice.sh and build_project.sh do; the
% lab packets build in ECE444-faculty from its vendored copy of this file.
%
% Compiling with pdflatex anyway is not an error: the guard below falls back
% to the LaTeX default so the source still compiles under any engine and can
% be syntax-checked. The resulting PDF just is not the one we ship.
%
% Paths are relative to latex/, which is where the build scripts run
% (each does cd "$(dirname "$0")"). Loading by path rather than by installed
% family name means no system font install and no fontconfig, so a fresh
% clone builds a byte-comparable PDF anywhere.
%
% MATH IS DELIBERATELY LEFT IN COMPUTER MODERN. Barlow has no math companion,
% and CM math under sans body text is already the house pairing: the decks set
% words in Source Sans Pro and let MathJax render symbols in its CM-derived
% face. Moving math to a sans as well would mean unicode-math plus a sans math
% font (Fira Math being the plausible one) and re-checking every equation in
% the course for symbol coverage. That is a separate decision, not a side
% effect of this one.

\usepackage{iftex}
\ifLuaTeX
	\usepackage{fontspec}
	\defaultfontfeatures{Ligatures=TeX}
	\setmainfont{Barlow}[
		Path           = fonts/Barlow/,
		Extension      = .ttf,
		UprightFont    = *-Regular,
		ItalicFont     = *-Italic,
		BoldFont       = *-Bold,
		BoldItalicFont = *-BoldItalic,
	]
	\setsansfont{Barlow}[
		Path           = fonts/Barlow/,
		Extension      = .ttf,
		UprightFont    = *-Regular,
		ItalicFont     = *-Italic,
		BoldFont       = *-SemiBold,
		BoldItalicFont = *-BoldItalic,
	]
	% microtype earns its place here specifically because of Barlow. Its
	% glyphs are narrower than Computer Modern's, so justified prose sets
	% tighter and TeX has to squeeze interword space to fit; two paragraphs
	% in the L02/L03 sets went overfull on the first Barlow build. Protrusion
	% and font expansion absorb that -- both are lualatex features, so this
	% belongs inside the guard with everything else.
	\usepackage{microtype}

	% Mono stays Latin Modern: Barlow ships no monospaced face, and verbatim
	% and code need one that actually is fixed-pitch.
\fi

\printselectedfalse      % full worked solutions in the key, not answers-only
\CorrectChoiceEmphasis{\color{red}}
\SolutionEmphasis{\color{red}}
\renewcommand\questionlabel{\textbf{\thequestion.}}
\renewcommand{\thepartno}{\alph{partno}}
\renewcommand{\partlabel}{(\thepartno)}

% Fillable table cell: blank in the student copy, the expected value in red on
% the KEY. Fixed-width centered columns keep the blank cells writable.
\newcolumntype{K}[1]{>{\centering\arraybackslash}p{#1}}
\newcommand{\kv}[1]{\ifkey{\color{red}#1}\fi\rule[-0.10in]{0pt}{0.34in}}
% A fixed (printed) cell, on the same baseline as the fillable ones.
\newcommand{\fv}[1]{#1\rule[-0.10in]{0pt}{0.34in}}

% print the answers if this is the key
\ifkey
	\printanswers % if this IS the key
\else
	\noprintanswers  %if this is NOT the key
\fi

% print the header and footer
% need to print key on top if this is the key
\ifkey
	\firstpageheader{}{\color{red}{ KEY}}{}
	\runningheader{\thisclass\hspace{1pt} -- \thisterm }{L\labLesson\ Lab\hspace{1pt}\color{red}{ KEY}}{\thepage\hspace{1pt} of\hspace{1pt} \numpages}
	\firstpagefooter{}{}{}
	\runningfooter{}{}{}
\else
	\firstpageheader{}{}{}
	\runningheader{\thisclass \hspace{1pt} -- \thisterm }{L\labLesson\ Lab\hspace{1pt}}{\thepage\hspace{1pt} of\hspace{1pt} \numpages}
	\firstpagefooter{}{}{}
	\runningfooter{}{}{}
\fi

\runningheadrule

\begin{document}
\pagestyle{headandfoot}   % exam class


%%%%%%%%%%%%%%%%%%%%%%%%%%%%%%%%%%%%%%%%%%%%%%%%%%%%%%
% title block
%%%%%%%%%%%%%%%%%%%%%%%%%%%%%%%%%%%%%%%%%%%%%%%%%%%%%%

\begin{center}
USAF ACADEMY \\
DEPARTMENT OF ELECTRICAL AND COMPUTER ENGINEERING \\
LESSON \labLesson\ LAB - Array Factor \\
\textbf{Lab Packet} \\
\thisterm \\[4mm]
\end{center}

\vspace{-4pt}
\textbf{Name:} \rule[-2pt]{2.6in}{0.4pt} \hfill
\textbf{Section:} \rule[-2pt]{0.9in}{0.4pt} \hfill
\textbf{Date:} \rule[-2pt]{1.1in}{0.4pt}

\vspace{4pt}

\fullwidth{\textbf{Learning Objective 3.2: } I can derive the array factor for an arbitrary linear
array and apply pattern multiplication.}

\vspace{2pt}

This packet is the L\labLesson\ Lab, in three parts, and you turn it in once,
complete, before Lesson 21. Do \textbf{Part 1} on your own before the lab
period: derive the array factor for each configuration you will sweep and
predict what the trace will show. You need it at the bench. \textbf{Part 2} is
your own measurements, from your own kit, recorded in class. Do \textbf{Part 3}
on your own after the lab: reconcile the measurements with the predictions and
work three analysis problems.

\textbf{Grading.} One submission, counted toward your engagement grade and
graded as a genuine, documented attempt.

The array is the PHASER's: $N = 8$ elements, E1 to E8, on a $d = 14\ \text{mm}$
pitch. Use $c = 3.00\times 10^{8}\ \m/\text{s}$.

\begin{questions}

\fullwidth{\large\textbf{Part 1: Pre-Lab} \normalsize(individual, before the lab period)}

%%%%%%%%%%%%%%%%%%%%%%%%%%%%%%%%%%%%%%%%%%%%%%%%%%%
\question \textbf{Your kit's frequency.} Every prediction in this packet depends
on the source frequency. Use the HB100 frequency your kit reported in the L18
Lab. If you did not record one, use the nominal $10.525\ \ghz$ and say so. At the
bench, recompute if \textbf{Find HB100} reports a frequency more than
$50\ \text{MHz}$ from the one you used here.

\begin{center}
Signal Freq = \ansbox[1.1in]{$10.3\ \ghz$} \hspace{0.25in}
$\lambda$ = \ansbox[1.1in]{$29.1$ mm} \\[10pt]
$d/\lambda$ = \ansbox[1.1in]{$0.481$} \hspace{0.25in}
$kd$ = \ansbox[1.1in]{$173.0\mydeg$}
\end{center}

\begin{solution}
This key works one nominal frequency, $10.3\ \ghz$, the one the L20 Lab page
uses for its worked example. Accept each cadet's own Signal Freq, anywhere in
the HB100's $10.1$ to $10.7\ \ghz$ range, and check the arithmetic from it.
\eq{ \lambda &= \frac{c}{f} = \frac{3.00\times 10^{8}}{10.3\times 10^{9}} = 29.1\ \text{mm} \\
     \frac{d}{\lambda} &= \frac{14}{29.1} = 0.481 \\
     kd &= 360\mydeg\ \frac{d}{\lambda} = 173.0\mydeg }
The dry-run kit (phaser-00, 2026-10-19) read $10.408\ \ghz$, where
$\lambda = 28.82$ mm, $d/\lambda = 0.486$, and $kd = 174.9\mydeg$.
\end{solution}

%%%%%%%%%%%%%%%%%%%%%%%%%%%%%%%%%%%%%%%%%%%%%%%%%%%
\question \textbf{Uniform apertures.} In Part 2 you sweep three contiguous,
centered apertures with every active gain at 100: all eight elements, E3 to E6,
and E4 and E5.

\begin{parts}

	\part Fill the prediction table at your frequency. State which method you used
	for each half-power beamwidth. The change from $N = 8$ is the ideal voltage-sum
	prediction for the peak.

	\begin{center}
	\small
	\renewcommand{\arraystretch}{1.25}
	\begin{tabular}{| l | K{0.65in} | K{0.75in} | K{0.75in} | K{0.75in} | K{0.95in} |}
	\hline
	\textbf{Active} & \textbf{$Nd$} & \textbf{HPBW} & \textbf{FNBW} & \textbf{First SLL} & \textbf{Change from $N = 8$} \\\hline\hline
	8 (E1--E8) & \kv{$112$ mm} & \kv{$13.3\mydeg$} & \kv{$30.1\mydeg$} & \kv{$-12.8$ dB} & \fv{reference} \\\hline
	4 (E3--E6) & \kv{$56$ mm} & \kv{$27.4\mydeg$} & \kv{$62.6\mydeg$} & \kv{$-11.3$ dB} & \kv{$-6.0$ dB} \\\hline
	2 (E4, E5) & \kv{$28$ mm} & \kv{$62.6\mydeg$} & \kv{$180\mydeg$} & \kv{none} & \kv{$-12.0$ dB} \\\hline
	\end{tabular}
	\end{center}

	\begin{solution}
	Keyed at $10.3\ \ghz$ from the exact array factor. The $0.886\lambda/(Nd)$
	rule gives $13.2\mydeg$ and $26.4\mydeg$ for the first two rows, and either is
	acceptable there. FNBW is $2\arcsin(\lambda/Nd)$. The peak change is
	$20\mylog{N_{\text{active}}/8}$. At $10.408\ \ghz$ the same columns read
	$13.2\mydeg$, $27.1\mydeg$, $61.9\mydeg$ and $29.8\mydeg$, $62.0\mydeg$.
	\end{solution}

	\begin{minipage}[t]{0.58\linewidth}
		\part For the two-element row, solve the half-power beamwidth directly from
		$AF = \cosp{\psi/2}$, $\psi = kd\sinp{\theta}$. Then evaluate the
		$0.886\lambda/(Nd)$ rule and state its percentage error.
		\begin{solution}[1.6in]
		\eq{ \cosp{\psi/2} &= \frac{1}{\sqrt{2}} \ \Rightarrow\ \psi = \frac{\pi}{2} \\
		     \sinp{\theta} &= \frac{\lambda}{4d} = \frac{29.1}{56} = 0.520 \\
		     \theta_{\text{HP}} &= 2(31.3\mydeg) = 62.6\mydeg }
		\eq{ \frac{0.886(29.1)}{28} &= 0.921\ \text{rad} = 52.8\mydeg \\
		     \text{error} &= \frac{52.8 - 62.6}{62.6} = -15.7\% }
		The rule assumes $\sinp{\theta} \approx \theta$ at the half-power point,
		and at $31\mydeg$ that no longer holds.
		\end{solution}
	\end{minipage}\hfill%
	\begin{minipage}[t]{0.37\linewidth}
		\raggedleft \vspace{12pt}
		$\theta_{\text{HP}}$ = \ansbox[1.3in]{$62.6\mydeg$}
		\\[10pt]
		error = \ansbox[1.3in]{$-15.7\%$}
	\end{minipage}

	\part Evaluate $\lambda/2d$ at your frequency. Use it to explain what you
	entered for the two-element FNBW and first SLL, and why.
		\begin{solution}[1.6in]
		The first null of $\cosp{\psi/2}$ needs $\psi = \pi$, which is
		$\sinp{\theta} = \lambda/2d = 29.1/28 = 1.04$. No real angle has a sine
		above 1, so the null falls outside visible space and the pattern has no
		null anywhere in the scan. The beamwidth between first nulls is quoted as
		the full $180\mydeg$ by convention. With no null there is no sidelobe
		either: $\cosp{\psi/2}$ falls monotonically from the peak to the edge of
		visible space, so the first SLL is ``none.'' The condition
		$\lambda/2d > 1$ holds below $c/2d = 10.71\ \ghz$, so every HB100 gives
		the same answer.
		\end{solution}

\end{parts}

%%%%%%%%%%%%%%%%%%%%%%%%%%%%%%%%%%%%%%%%%%%%%%%%%%%
\newpage
\question \textbf{Two elements, far apart.} The last sweep in Part 2 leaves only
E1 and E8 on, so the array is two elements again, but now $D = 7d = 98\ \text{mm}$
apart instead of $14\ \text{mm}$.

\begin{parts}

	\part Write the array factor for this pair and say how it differs from the
	E4 and E5 pair.
		\begin{solution}[1.0in]
		It is the same two-element array factor with $D$ in place of $d$:
		\eq{ AF = \cosp{\frac{\psi}{2}}, \qquad \psi = kD\sinp{\theta}, \qquad
		     \frac{D}{\lambda} = \frac{98}{29.1} = 3.37 }
		The spacing is $7$ times larger, so $\psi$ runs through $7$ times as many
		periods across the scan.
		\end{solution}

	\begin{minipage}[t]{0.58\linewidth}
		\part Find every angle in visible space where the array factor reaches its
		full peak.
		\begin{solution}[1.5in]
		Full peaks sit at $\psi = 2\pi m$, which is
		\eq{ \sinp{\theta_m} = \frac{m\lambda}{D} = 0.297\,m }
		so $m = 0, \pm 1, \pm 2, \pm 3$ give $0\mydeg$, $\pm 17.3\mydeg$,
		$\pm 36.5\mydeg$, and $\pm 63.1\mydeg$. At $m = 4$ the sine is $1.19$,
		outside visible space. All seven lobes are full height: these are grating
		lobes, and Lesson 23 takes them up.
		\end{solution}
	\end{minipage}\hfill%
	\begin{minipage}[t]{0.37\linewidth}
		\raggedleft \vspace{12pt}
		peaks = \ansbox[1.5in]{$0, \pm 17.3, \pm 36.5, \pm 63.1\mydeg$}
	\end{minipage}

	\begin{minipage}[t]{0.58\linewidth}
		\part Find the half-power beamwidth and the first-null beamwidth of the lobe
		at $0\mydeg$.
		\begin{solution}[1.4in]
		\eq{ \sinp{\theta} &= \frac{\lambda}{4D} = 0.0743 \ \Rightarrow\
		     \theta_{\text{HP}} = 2(4.26\mydeg) = 8.5\mydeg \\
		     \sinp{\theta} &= \frac{\lambda}{2D} = 0.149 \ \Rightarrow\
		     \text{FNBW} = 2(8.54\mydeg) = 17.1\mydeg }
		\end{solution}
	\end{minipage}\hfill%
	\begin{minipage}[t]{0.37\linewidth}
		\raggedleft \vspace{12pt}
		$\theta_{\text{HP}}$ = \ansbox[1.3in]{$8.5\mydeg$}
		\\[10pt]
		FNBW = \ansbox[1.3in]{$17.1\mydeg$}
	\end{minipage}

	\part Predict the change from $N = 8$ for this pair's peak. Then explain, in two
	or three sentences, why two elements give a main lobe narrower than all eight
	do, and why the trace shows several lobes instead of one.
		\begin{solution}[1.6in]
		Two elements at gain 100 add to twice one element's voltage wherever they
		sit, so the ideal change is the same as for the E4 and E5 pair:
		\eq{ 20\mylog{\frac{2}{8}} = -12.0\ \db } The width of a lobe is set by how far apart the outermost
		elements are in wavelengths, and E1 and E8 span $98$ mm, nearly the full
		array, so the lobe is as narrow as a full aperture's. With only two
		samples of that aperture, spaced more than a wavelength apart, the phase
		difference returns to a multiple of $2\pi$ at several real angles, and
		the main lobe repeats at each one.
		\end{solution}

\end{parts}

%%%%%%%%%%%%%%%%%%%%%%%%%%%%%%%%%%%%%%%%%%%%%%%%%%%
\newpage
\question \textbf{Simulator rehearsal.} Open the hosted simulator at
\url{https://livethisdream.github.io/phaser/}. Its source is fixed at boresight,
so it can rehearse every sweep in Part 2.

\begin{enumerate}
	\item Open the \textbf{Lab Presets} section, press \textbf{2 Array Factor},
	then press \textbf{Start}. Record the Signal Freq it reports.
	\item In \textbf{Element Gains}, set E2 through E7 to 0, leaving E1 and E8 at
	100. Press \textbf{Start}.
	\item On the \textbf{Rectangular} tab, read the angles of the full-height lobes
	and the half-power width of the lobe at $0\mydeg$.
\end{enumerate}

\begin{center}
\renewcommand{\arraystretch}{1.25}
\begin{tabular}{| K{1.3in} | K{2.4in} | K{1.3in} |}
\hline
\textbf{Signal Freq} & \textbf{Full-height lobes} & \textbf{HPBW at $0\mydeg$} \\\hline\hline
\kv{$10.525\ \ghz$} & \kv{$0, \pm 16.9, \pm 35.6, \pm 60.8\mydeg$} & \kv{$8.3\mydeg$} \\\hline
\end{tabular}
\end{center}

\begin{solution}
Keyed from the array factor at the simulator's $10.525\ \ghz$,
$\lambda = 28.5$ mm, $D/\lambda = 3.44$. The simulator adds the element voltages
exactly, so its lobes should land within one sweep step, about $1\mydeg$ near
broadside and $2\mydeg$ near $60\mydeg$, of these angles. A cadet who rehearsed
at $10.525\ \ghz$ and measures at their own frequency will see the outer lobes
move by a degree or two; that is the frequency, not an error.
\end{solution}

In one sentence, compare the simulator's lobe angles with your Question 3(b), and
say why any difference is expected.

\begin{solution}[0.9in]
They agree once both are at the same frequency; the lobe angles scale with
$\lambda/D$, so a prediction made at another frequency shifts the outer lobes.
\end{solution}

%%%%%%%%%%%%%%%%%%%%%%%%%%%%%%%%%%%%%%%%%%%%%%%%%%%
\newpage
\fullwidth{\large\textbf{Part 2: Measurements} \normalsize(your kit, in class)}

Follow the procedure on the L20 Lab page in read mode. Bring the kit up as in the
L17 Lab: the HB100 at boresight, about $1\ \m$ out at patch-row height, then
\textbf{Find HB100}, then \textbf{Calibrate}. Load \textbf{2 Array Factor} under
\textbf{Lab Presets}, and set \textbf{Rx Gain (dB)} for an FFT peak of
$-10 \pm 3$ dBFS. Leave the gain there for all four sweeps, so the peaks compare
against each other. Read each sidelobe or lobe level relative to that trace's own
peak.

\vspace{2pt}
\begin{center}
Signal Freq after Find HB100 = \ansbox[1.3in]{your kit's} \hspace{0.25in}
Rx Gain = \ansbox[1.1in]{as set}
\end{center}

%%%%%%%%%%%%%%%%%%%%%%%%%%%%%%%%%%%%%%%%%%%%%%%%%%%
\question \textbf{Table A: four apertures.} One row per sweep.

\begin{center}
\small
\setlength{\tabcolsep}{3pt}
\renewcommand{\arraystretch}{1.25}
\begin{tabular}{| l | K{0.8in} | K{0.9in} | K{0.8in} | K{0.8in} | K{1.1in} |}
\hline
\textbf{Active} & \textbf{Peak (dBFS)} & \textbf{Change from $N = 8$ (dB)} & \textbf{HPBW} & \textbf{FNBW} & \textbf{First SLL (dBc)} \\\hline\hline
8 (E1--E8) & \kv{$-5.25$} & \fv{reference} & \kv{$13.4\mydeg$} & \kv{$29.8\mydeg$} & \kv{$-12.5$ to $-14.0$} \\\hline
4 (E3--E6) & \kv{$-10.4$} & \kv{$-5.1$} & \kv{$27.4\mydeg$} & \kv{$\approx 67\mydeg$} & \kv{$-12$ to $-14$} \\\hline
2 (E4, E5) & \kv{$-15.1$} & \kv{$-9.9$} & \kv{$\approx 56\mydeg$} & \kv{no null} & \kv{none} \\\hline
2 (E1, E8) & \kv{} & \kv{near the E4, E5 row} & \kv{$8.5\mydeg \pm 1\mydeg$} & \kv{$17\mydeg \pm 1\mydeg$} & \kv{lobes near $0$} \\\hline
\end{tabular}
\end{center}

\begin{solution}
The first three rows are the dry run's, on phaser-00 at $10.408\ \ghz$, Rx Gain
$1$ dB, source about $1.2$ m out. The E1 and E8 row was not measured in the dry
run: expect the main lobe within about $1\mydeg$ of Question 3(c), and a peak
within about $1$ dB of the E4 and E5 row, since both rows add two equal voltages.
For that row the last column holds the level of the lobes at $\pm 17\mydeg$, which
the array factor puts at full height. Absolute dBFS values depend on the Rx Gain
the cadet set and are not graded.
\end{solution}

%%%%%%%%%%%%%%%%%%%%%%%%%%%%%%%%%%%%%%%%%%%%%%%%%%%
\question \textbf{Table B: the E1 and E8 lobes.} From the E1 and E8 trace, read
the angle and the level (relative to the trace's peak) of each full-height lobe.

\begin{center}
\small
\renewcommand{\arraystretch}{1.25}
\begin{tabular}{| l | K{0.58in} | K{0.58in} | K{0.58in} | K{0.58in} | K{0.58in} | K{0.58in} | K{0.58in} |}
\hline
\textbf{Lobe $m$} & $-3$ & $-2$ & $-1$ & $0$ & $+1$ & $+2$ & $+3$ \\\hline\hline
Angle & \kv{$-63\mydeg$} & \kv{$-36\mydeg$} & \kv{$-17\mydeg$} & \kv{$0\mydeg$} & \kv{$17\mydeg$} & \kv{$36\mydeg$} & \kv{$63\mydeg$} \\\hline
Level (dBc) & \kv{$\approx 0$} & \kv{$\approx 0$} & \kv{$\approx 0$} & \fv{$0$} & \kv{$\approx 0$} & \kv{$\approx 0$} & \kv{$\approx 0$} \\\hline
\end{tabular}
\end{center}

\begin{solution}
Keyed from the array factor; each angle within one sweep step of Question 3(b)
at the cadet's frequency, and the outer lobes widest. Accept levels within a dB
or two of the main lobe: the sweep holds the element pattern fixed at the
source's direction, so the grating lobes appear at full height on the trace.
\end{solution}

%%%%%%%%%%%%%%%%%%%%%%%%%%%%%%%%%%%%%%%%%%%%%%%%%%%
\newpage
\fullwidth{\large\textbf{Part 3: Writeup} \normalsize(individual, after the lab)}

%%%%%%%%%%%%%%%%%%%%%%%%%%%%%%%%%%%%%%%%%%%%%%%%%%%
\question \textbf{Reconcile the measurements.} Answer each part in three or four
sentences, using your numbers from Part 2 and your own predictions from Part 1.

\begin{parts}

	\part Compare the measured beamwidths of the three uniform apertures with
	Question 2. Does the beam double at each halving? Name the cause of any
	difference larger than a degree.
	\begin{solution}[1.5in]
	The eight- and four-element HPBW should land within about $1\mydeg$ of the
	calculated values, and the beam doubles at each halving because $Nd$ halves.
	The two-element HPBW reads narrow and scattered (the dry run read $52$ to
	$58\mydeg$ against $61.9\mydeg$), because that pattern falls only about
	$0.2$ dB per degree at its half-power points, so a small error in the level
	moves each crossing by degrees. The four-element FNBW reads wide when one
	null fills to the floor, about $-23$ dBc. Sidelobes are a range, not a single
	number.
	\end{solution}

	\part Compare the E1 and E8 row with the E4 and E5 row. Both have two active
	elements. Explain, with your numbers, why their peaks agree and their beams do
	not.
	\begin{solution}[1.5in]
	The peaks agree within about a dB because each is the sum of two equal
	voltages, which does not depend on where the elements sit. The beams differ
	because the width of a lobe is set by the span of the elements in
	wavelengths: $28$ mm for E4 and E5, $98$ mm for E1 and E8. The wide pair gives
	a main lobe of about $8.5\mydeg$, narrower than the full eight-element beam,
	but with spacing above a wavelength it repeats that lobe at full height near
	$\pm 17$, $\pm 36$, and $\pm 63\mydeg$.
	\end{solution}

	\part The voltage sum predicts a $6.0$ dB peak change per halving. Compare your
	two measured changes with it. Leakage through the elements set to 0 and
	receiver compression were ruled out on the dry-run kit; propose one other
	possible cause of a shortfall, and describe a measurement on your kit that
	would test it.
	\begin{solution}[2.0in]
	The dry-run kit read $5.1$ and $4.8$ dB; expect about $5 \pm 1$ dB per
	halving. The cause is not yet known, so grade the proposal, not a match to an
	answer. Give credit for any cause the cadet ties to a measurement that could
	confirm or rule it out. Two candidates are still open, and either is a good
	answer: Peak Array Gain takes the highest of 162 noisy samples, so a flat,
	small-aperture trace reads high (test: compare the median of one element's
	sweep with its maximum); or channel phase or gain shifts with the gain
	setting, so the signals stop adding fully in phase away from the calibration
	setting (test: measure each channel's phase at gains 100 and 50).
	\end{solution}

\end{parts}

%%%%%%%%%%%%%%%%%%%%%%%%%%%%%%%%%%%%%%%%%%%%%%%%%%%
\newpage
\question \textbf{Reading a sweep.} A cadet sweeps the PHASER with all eight
elements active and a source at boresight, at $10.3\ \ghz$. The table lists every
third sample of the trace, about $2.8\mydeg$ apart. The peak reads $-4.1$ dBFS.

\begin{center}
\begin{tabular}{r|rrrrrrrrrr}
$\theta$ (\mydeg) & $0.0$ & $2.8$ & $5.6$ & $8.4$ & $11.3$ & $14.1$ & $16.9$ & $19.7$ & $22.5$ & $25.3$ \\ \hline
Power (dBFS) & $-4.1$ & $-4.6$ & $-6.2$ & $-9.1$ & $-14.3$ & $-24.1$ & $-22.0$ & $-17.4$ & $-16.6$ & $-18.1$ \\
\end{tabular}
\end{center}

\begin{parts}

	\begin{minipage}[t]{0.58\linewidth}
		\part Using linear interpolation between the two listed samples that bracket
		the half-power level, find the half-power beamwidth. The trace is symmetric
		about $0\mydeg$.
		\begin{solution}[1.3in]
		The half-power level is $-4.1 - 3.0 = -7.1$ dBFS, between $5.6\mydeg$
		($-6.2$) and $8.4\mydeg$ ($-9.1$).
		\eq{ \theta_{-3} &= 5.6 + 2.8\lp\frac{7.1-6.2}{9.1-6.2}\rp = 6.5\mydeg \\
		     \theta_{\text{HP}} &= 2(6.5) = 12.9\mydeg }
		Theory gives $13.3\mydeg$, so the coarse table reads a few tenths narrow.
		\end{solution}
	\end{minipage}\hfill%
	\begin{minipage}[t]{0.37\linewidth}
		\raggedleft \vspace{12pt}
		$\theta_{\text{HP}}$ = \ansbox[1.3in]{$12.9\mydeg$}
	\end{minipage}

	\begin{minipage}[t]{0.58\linewidth}
		\part Find the first-null beamwidth from the listed samples, and compare it
		with the array factor's value. Say why they differ.
		\begin{solution}[1.4in]
		The deepest listed sample is $-24.1$ dBFS at $14.1\mydeg$, so the table
		gives $2(14.1) = 28.2\mydeg$, against
		\eq{ \text{FNBW} = 2\arcsin\lp\frac{29.1}{112}\rp = 30.1\mydeg }
		The true null at $15.1\mydeg$ falls between listed samples, and the table
		can only report a null at an angle it lists. The full sweep, one sample
		every $0.93\mydeg$, lands within a step of it.
		\end{solution}
	\end{minipage}\hfill%
	\begin{minipage}[t]{0.37\linewidth}
		\raggedleft \vspace{12pt}
		FNBW = \ansbox[1.3in]{$28.2\mydeg$}
	\end{minipage}

	\part No sample reads below about $-24$ dBFS, $20$ dB below the peak, yet the
	floor on this kit is $23$ dB below the peak. Explain in one or two sentences
	why the deepest sample sits $3$ dB above the floor rather than on it.
		\begin{solution}[1.1in]
		The receiver reports signal power plus noise power. Near the null the array
		factor has fallen to about the floor level, so the two contributions are
		equal and their sum is $3$ dB above either one.
		\end{solution}

\end{parts}

%%%%%%%%%%%%%%%%%%%%%%%%%%%%%%%%%%%%%%%%%%%%%%%%%%%
\question \textbf{Reconciling a sidelobe.} A cadet measures the first sidelobe of
the full eight-element array at $-11.5$ dBc, while the array factor for $N = 8$
predicts $-12.8$ dB. The noise floor sits at $-23$ dBc.

\begin{parts}

	\begin{minipage}[t]{0.58\linewidth}
		\part Treating the floor and the sidelobe as uncorrelated powers, compute
		the level the sweep should report for a true $-12.8$ dB sidelobe on a
		$-23$ dBc floor.
		\begin{solution}[1.3in]
		\eq{ p &= 10^{-12.8/10} + 10^{-23/10} \\
		     &= 0.0525 + 0.0050 = 0.0575 \\
		     \text{SLL} &= 10\mylog{0.0575} = -12.4\ \text{dBc} }
		The floor accounts for about $0.4$ dB of lift; it lifts the peak by only
		$0.02$ dB.
		\end{solution}
	\end{minipage}\hfill%
	\begin{minipage}[t]{0.37\linewidth}
		\raggedleft \vspace{12pt}
		SLL = \ansbox[1.3in]{$-12.4$ dBc}
	\end{minipage}

	\part The floor explains $0.4$ dB of the $1.3$ dB difference. Name two other
	contributions that plausibly make up the rest, and say for each whether it
	moves the reading up or down.
		\begin{solution}[1.4in]
		Run-to-run noise on the sidelobe sample, which sits only about $10$ dB
		above the floor and scatters a few tenths of a decibel either way; and
		residual amplitude and phase errors across the elements after
		calibration, which break the uniform distribution and raise the
		sidelobes. Grid straddle moves the reading down, and by very little here:
		the sidelobe peak at $21.9\mydeg$ falls within a hundredth of a decibel of
		a sample.
		\end{solution}

\end{parts}

%%%%%%%%%%%%%%%%%%%%%%%%%%%%%%%%%%%%%%%%%%%%%%%%%%%
\question \textbf{Design judgment.} You are asked to demonstrate a $20\mydeg$ beam
on the PHASER at $10.3\ \ghz$, using only the Element Gains controls.

\begin{parts}

	\begin{minipage}[t]{0.58\linewidth}
		\part How many adjacent elements must be active to give a half-power
		beamwidth closest to $20\mydeg$? Give the aperture and the resulting
		beamwidth.
		\begin{solution}[1.4in]
		\eq{ Nd &= \frac{0.886\lambda}{\theta_{\text{HP}}} = \frac{0.886(29.1)}{0.349} = 73.8\ \text{mm} \\
		     N &= \frac{73.8}{14} = 5.3 }
		Take $N = 5$: $Nd = 70$ mm and
		$\theta_{\text{HP}} \approx 0.886(29.1)/70 = 0.368\ \text{rad} = 21.1\mydeg$.
		Six elements give $17.6\mydeg$, farther from the target.
		\end{solution}
	\end{minipage}\hfill%
	\begin{minipage}[t]{0.37\linewidth}
		\raggedleft \vspace{12pt}
		$N$ = \ansbox[1.3in]{$5$}
		\\[10pt]
		$\theta_{\text{HP}}$ = \ansbox[1.3in]{$21.1\mydeg$}
	\end{minipage}

	\part A classmate proposes E1 and E8 alone instead, since a two-element array
	``has a wide beam.'' Using your Part 2 data, explain in two sentences why that
	does not produce a $20\mydeg$ beam.
		\begin{solution}[1.1in]
		The lobe width is set by the span of the elements, and E1 and E8 span
		$98$ mm, so their main lobe is about $8.5\mydeg$, not wide. The pair also
		repeats that lobe at full height every $17\mydeg$ or so, which is the
		opposite of one $20\mydeg$ beam.
		\end{solution}

	\part Another classmate suggests a Hann taper on all eight elements, which
	widens the beam to about $19.5\mydeg$. State one advantage and one disadvantage
	of the taper against switching elements off, and say which you would choose for
	a demonstration whose purpose is to show that beamwidth depends on aperture.
		\begin{solution}[1.5in]
		The taper keeps the full $112$ mm aperture and pushes the sidelobes below
		the floor, so the trace is cleaner. But it widens the beam by changing the
		shape of the amplitude distribution, not the aperture, so it demonstrates
		the opposite of the intended point, and it loses about $4.7$ dB of peak.
		For a demonstration about aperture, switch elements off: the beamwidth
		and the aperture then move together in one control.
		\end{solution}

\end{parts}

\end{questions}

\vspace{1.0em}
\noindent\textbf{Documentation:} \rule[-2pt]{0.6\linewidth}{0.4pt}

\end{document}
